% !PRINTED 353
\weberpage{353}
% [sic] The source introduces \\mathfrak{m} but predicates about \\mathfrak{a} before defining it; the wording is retained.
%<<CONT

2. Theorem. \gesperrt{Every number of the order $[Q]$, and only these, are contained in the form}
\weq{6}{\omega = x + y\theta,}
\gesperrt{in which $x, y$ are arbitrary rational integers.}

We call the system
\weq{7}{(1, \theta)}
a basis of the order $[Q]$.

By substituting (5) one can also represent the numbers (6) in the form:
\weq{8}{\omega = \dfrac{x + y\sqrt{D}}{2},}
where the $x, y$ are integers which satisfy the condition
\weq{9}{x^2 - y^2 D \equiv 0 \pmod{4}.}
The number $D$ is a discriminant whose associated fundamental discriminant is $\Delta$. We call it the \gesperrt{discriminant of the order $[Q]$}.

\begin{center}
\S\ 97. \gesperrt{Orders and Ideals.}
\end{center}

Let $\mathfrak{m}$ be an ideal of the quadratic field $\Omega$. If all integers divisible by $\mathfrak{a}$ are divisible by a rational number $m$, then the ideal $\mathfrak{a} = m\mathfrak{m}$ is divisible by $m$; we have
\[ N(\mathfrak{a}) = m^2 N(\mathfrak{m}), \]
and the rational number $m N(\mathfrak{m})$, which is less than $N(\mathfrak{a})$, is divisible by $\mathfrak{a}$.

If $m$ is the greatest natural number by which $\mathfrak{a}$ is divisible, we shall call $m$ \gesperrt{the divisor of $\mathfrak{a}$}. If the divisor $= 1$, the ideal is called \gesperrt{primary}, otherwise \gesperrt{derived}. One obtains all ideals derived from a primary ideal $\mathfrak{m}$ by multiplying $\mathfrak{m}$ by arbitrary natural numbers, and all ideals derived from a primary ideal are \gesperrt{equivalent} to one another.

3. Theorem. \gesperrt{If $\mathfrak{a}$ is a primary ideal, then the norm of $\mathfrak{a}$ is at the same time the least natural number divisible by $\mathfrak{a}$.}

For let us denote by $a$ the least natural number divisible by $\mathfrak{a}$; the norm of $\mathfrak{a}$ is then in any case divisible by $a$. We put
\[ N(\mathfrak{a}) = ma, \]
and if $\mathfrak{a}'$ is the ideal conjugate to $\mathfrak{a}$, then, since by Bd.\ II, \S\ 164 the norm of an ideal is equal to the product of the conjugate ideals, it follows that
\weq{1}{N(\mathfrak{a}) = \mathfrak{a}\mathfrak{a}' = ma.}
